%!TEX program = xelatex
% ============================================================
%  הרצאה 3 — תנע, שימור התנע, וכוחות בטבע
% ============================================================
\documentclass[12pt, a4paper]{article}

\usepackage{amsmath, amssymb, mathtools}
\usepackage{xcolor}
\usepackage{graphicx}

\usepackage{tcolorbox}
\tcbuselibrary{skins, breakable}
\tcbset{
  resultbox/.style={
    enhanced, colback=white, colframe=black,
    boxrule=0.8pt, arc=3pt,
    left=2pt, right=2pt, top=1.5pt, bottom=1.5pt,
  },
  notebox/.style={
    enhanced, breakable,
    colback=gray!8, colframe=gray!50,
    boxrule=0.8pt, arc=3pt,
    left=2pt, right=2pt, top=1.5pt, bottom=1.5pt,
    fonttitle=\bfseries,
  },
  warnbox/.style={
    enhanced, breakable,
    colback=red!5, colframe=red!60,
    boxrule=0.8pt, arc=3pt,
    left=3pt, right=3pt, top=2pt, bottom=2pt,
    fonttitle=\bfseries,
  },
  defbox/.style={
    enhanced, breakable,
    colback=blue!5, colframe=blue!50,
    boxrule=0.8pt, arc=3pt,
    left=3pt, right=3pt, top=2pt, bottom=2pt,
    fonttitle=\bfseries,
  },
}

\usepackage{tikz}
\usetikzlibrary{arrows.meta, calc, decorations.pathreplacing, patterns}
\usepackage[a4paper, top=2cm, bottom=2cm,
            left=2cm, right=2cm]{geometry}
\usepackage{setspace}
\setstretch{1.3}
\usepackage{titlesec}
\titleformat{\section}{\large\bfseries}{\thesection.}{0.5em}{}[\titlerule]
\titleformat{\subsection}{\bfseries}{\thesubsection}{0.5em}{}

\setlength{\headheight}{15pt}
\usepackage{fancyhdr}
\pagestyle{fancy}
\fancyhf{}
\rhead{\textenglish{\textbf{Lecture 3 --- Momentum \& Forces}}}
\lhead{\thepage}
\renewcommand{\headrulewidth}{0.3pt}

\usepackage{fontspec}
\usepackage{polyglossia}
\setmainlanguage{hebrew}
\setotherlanguage{english}
\setmainfont{Times New Roman}
\newfontfamily\hebrewfont{Times New Roman}
\newfontfamily\englishfont{Times New Roman}
\newfontfamily\hebrewfonttt{Courier New}

\newcommand{\hvec}[1]{\overrightarrow{#1}}

% ============================================================
% Reusable stick figure macros (templates)
% ============================================================
% Child standing (passive pose, arms slightly out).
% Usage inside a tikzpicture: \childStanding{x}{y}{scale}
\newcommand{\childStanding}[3]{%
  \begin{scope}[shift={(#1, #2)}, scale=#3, line width=0.8pt]
    \draw[fill=orange!25] (0, 1.6) circle (0.22);                % head
    \fill (-0.08, 1.66) circle (0.022);                          % left eye
    \fill ( 0.08, 1.66) circle (0.022);                          % right eye
    \draw (-0.08, 1.52) .. controls (0, 1.48) .. (0.08, 1.52);   % smile
    \draw (0, 1.38) -- (0, 0.55);                                % body
    \draw (0, 1.20) -- (-0.28, 0.80);                            % left arm
    \draw (0, 1.20) -- ( 0.28, 0.80);                            % right arm
    \draw (0, 0.55) -- (-0.20, 0);                               % left leg
    \draw (0, 0.55) -- ( 0.20, 0);                               % right leg
  \end{scope}
}

% Child stick figure throwing a ball to the right.
% Usage inside a tikzpicture: \childThrowingRight{x}{y}{scale}
%   x, y  -- position of the figure's feet (between the legs)
%   scale -- size multiplier (1.0 is default; ~0.8-1.2 typical)
\newcommand{\childThrowingRight}[3]{%
  \begin{scope}[shift={(#1, #2)}, scale=#3, line width=0.8pt]
    \draw[fill=orange!25] (0, 1.6) circle (0.22);                % head
    \fill (-0.08, 1.66) circle (0.022);                          % left eye
    \fill ( 0.08, 1.66) circle (0.022);                          % right eye
    \draw (-0.08, 1.52) .. controls (0, 1.48) .. (0.08, 1.52);   % smile
    \draw (0, 1.38) -- (0, 0.55);                                % body
    \draw (0, 1.20) -- ( 0.65, 1.15);                            % throwing arm
    \draw (0, 1.20) -- (-0.22, 0.75);                            % other arm
    \draw (0, 0.55) -- (-0.20, 0);                               % left leg
    \draw (0, 0.55) -- ( 0.20, 0);                               % right leg
  \end{scope}
}
% For a left-throwing child, call inside [xscale=-1]:
%   \begin{scope}[xscale=-1] \childThrowingRight{-x}{y}{scale} \end{scope}

% ============================================================
\begin{document}
% ============================================================

% ── Title ────────────────────────────────────────────────────
\begin{center}
  {\LARGE\bfseries הרצאה 3 --- תנע, שימור התנע, וכוחות בטבע}\\[6pt]
  \rule{\linewidth}{0.8pt}
\end{center}
\vspace{0.5cm}

% ============================================================
\section*{1. הקשר בין מתקף לחוק השני: התנע}
% ============================================================

נחבר את שתי המסקנות מההרצאה הקודמת. נתחיל מחוק שני,
\[
  \hvec{F} = m\hvec{a}
\]
ונבצע אינטגרציה לפי הזמן בקטע \textenglish{$[t_1, t_2]$}:
\[
  \int_{t_1}^{t_2} \hvec{F}\,dt = \int_{t_1}^{t_2} m\hvec{a}\,dt
  = \int_{t_1}^{t_2} m\frac{d\hvec{v}}{dt}\,dt = m\int_{\hvec{v}(t_1)}^{\hvec{v}(t_2)} d\hvec{v}
\]

עבור מסה קבועה:

\begin{tcolorbox}[resultbox]
\[
  \hvec{I} = m\hvec{v}(t_2) - m\hvec{v}(t_1)
\]
\end{tcolorbox}

האגף הימני "מתבקש" להגדרה משלו:

\begin{tcolorbox}[defbox, title={תנע (\textenglish{Momentum})}]
\[
  \hvec{p} \equiv m\hvec{v}
\]
\textbf{כמות התנועה} של גוף --- גודל וקטורי, יחידות \textenglish{kg$\cdot$m/s}.
\end{tcolorbox}

ובניסוח החדש:

\begin{tcolorbox}[resultbox, title={משפט מתקף-תנע}]
\[
  \hvec{I} = \Delta\hvec{p} = \hvec{p}(t_2) - \hvec{p}(t_1)
\]
\end{tcolorbox}

\begin{tcolorbox}[notebox, title={מסקנה רעיונית}]
\textbf{מתקף הוא בדיוק השינוי בתנע}. הוא לא איזה גודל מומצא --- הוא פשוט "כמה תנע הכוח העביר לגוף".
\end{tcolorbox}

\subsection*{1.1 ניסוח חלופי של החוק השני}

נחלק את משוואת המתקף-תנע ב-\textenglish{$dt$} (גבול אינפיניטסימלי):

\begin{tcolorbox}[resultbox]
\[
  \hvec{F} = \frac{d\hvec{p}}{dt}
\]
\end{tcolorbox}

זוהי הצורה \textbf{המקורית} של החוק השני (כפי שניוטון ניסח אותו). היא כללית יותר מ-\textenglish{$\hvec{F} = m\hvec{a}$}, כי היא תקפה גם כאשר המסה משתנה (למשל רקטה הפולטת דלק).

% ============================================================
\section*{2. שימור התנע במערכת סגורה}
% ============================================================

נחשב את התנע הכולל של מערכת סגורה של שני גופים \textenglish{$A, B$}:
\[
  \hvec{P}_{\text{tot}} = \hvec{p}_A + \hvec{p}_B
\]

נגזור לפי הזמן ונשתמש בחוק השני:
\[
  \frac{d\hvec{P}_{\text{tot}}}{dt} = \hvec{F}_A + \hvec{F}_B
  = \hvec{F}_{B\to A} + \hvec{F}_{A\to B}
\]

לפי החוק השלישי, \textenglish{$\hvec{F}_{B\to A} = -\hvec{F}_{A\to B}$}, ולכן הסכום מתאפס:

\begin{tcolorbox}[resultbox, title={חוק שימור התנע}]
\[
  \frac{d\hvec{P}_{\text{tot}}}{dt} = 0
  \qquad\Longrightarrow\qquad
  \hvec{P}_{\text{tot}} = \text{const}
\]
\end{tcolorbox}

\begin{tcolorbox}[notebox]
תנע אינו נוצר ואינו נעלם במערכת סגורה --- הוא רק \textbf{עובר} בין גופים דרך כוחות פנימיים. החוק השלישי הוא בעצם החוק הזה במסווה.
\end{tcolorbox}

% ============================================================
\section*{3. דוגמה: תפיסת כדור על קרח}
% ============================================================

ילד במסה \textenglish{$M = 60$ kg} עומד על קרח (חיכוך מתבטל), במנוחה. כדור במסה \textenglish{$m = 0.5$ kg} טס לעברו במהירות \textenglish{$|\hvec{v}_0| = 5$ m/s}. הילד תופס את הכדור.

\begin{center}
\begin{tikzpicture}[scale=1.2, >={Stealth[length=2.5mm]}]
  % ice surface
  \draw[thick] (-4, 0) -- (4, 0);
  \fill[pattern=north east lines, pattern color=blue!30] (-4, 0) rectangle (4, -0.4);
  \node[blue!60!black] at (3.2, -0.7) {\small קרח};
  % Child standing (left side)
  \childStanding{-1.5}{0}{1.0}
  \node[left] at (-1.85, 0.85) {\textenglish{$M$}};
  % Ball (right side, incoming)
  \fill[orange] (2.2, 1.15) circle (0.14);
  \node[above] at (2.2, 1.35) {\textenglish{$m$}};
  % Ball velocity arrow (pointing LEFT - toward child)
  \draw[->, very thick, red] (2, 1.15) -- (0.5, 1.15);
  \node[red, above] at (1.25, 1.2) {\textenglish{$\hvec{v}_0$}};
  % After-catch velocity of child (small arrow, same direction as v0)
  \draw[->, very thick, green!55!black, dashed] (-1.5, 2.4) -- (-2.4, 2.4);
  \node[green!55!black, above] at (-1.95, 2.45) {\textenglish{$\hvec{u}$ (לאחר התפיסה)}};
\end{tikzpicture}
\end{center}

\subsection*{3.1 מהירות הילד+כדור לאחר התפיסה}

המערכת סגורה (אין כוח חיצוני אופקי --- קרח חלק).

\textbf{תנע התחלתי} (לפני התפיסה): רק הכדור נע, הילד במנוחה.
\[
  \hvec{p}_{\text{before}} = m\hvec{v}_0
\]

\textbf{תנע סופי} (אחרי התפיסה): הילד והכדור נעים יחד במהירות \textenglish{$\hvec{u}$}.
\[
  \hvec{p}_{\text{after}} = (m + M)\hvec{u}
\]

מהשימור:

\begin{tcolorbox}[resultbox]
\[
  m\hvec{v}_0 = (m + M)\hvec{u}
  \qquad\Longrightarrow\qquad
  \hvec{u} = \frac{m}{m + M}\hvec{v}_0
\]
\end{tcolorbox}

מספרית:
\[
  |\hvec{u}| = \frac{0.5 \cdot 5}{0.5 + 60} \approx 0.041\;\text{m/s}
\]
(הילד נע בכיוון הכדור, אבל לאט בהרבה.)

\subsection*{3.2 הכוח שהכדור הפעיל על הילד בזמן התפיסה}

נניח שתהליך התפיסה ארך \textenglish{$\Delta t = 0.2$} שניות. נשתמש במשפט מתקף-תנע על הילד:
\[
  \hvec{F} \cdot \Delta t = \Delta\hvec{p}_{\text{child}} = M\hvec{u} - 0
\]

\begin{tcolorbox}[resultbox]
\[
  |\hvec{F}| = \frac{M\,|\hvec{u}|}{\Delta t}
  = \frac{60 \cdot 0.041}{0.2}
  \approx 12.5\;\text{N}
\]
\end{tcolorbox}

לפי החוק השלישי, אותו כוח (בכיוון הפוך) פעל מהילד על הכדור, ובלם אותו ממהירות \textenglish{$\hvec{v}_0$} למהירות \textenglish{$\hvec{u}$}.

% ============================================================
\section*{4. כוחות בטבע --- מבוא}
% ============================================================

חוקי ניוטון נותנים את \textbf{מסגרת} הדינמיקה. כעת נכיר את הכוחות הספציפיים שמופיעים בבעיות. עיקרי הכוחות:
\begin{enumerate}
  \item כוח הכבידה
  \item מתח בחוט
  \item כוח הנורמל
  \item חיכוך (סטטי וקינטי)
  \item כוחות חשמליים, גרזה, וכו' (יילמדו בנפרד)
\end{enumerate}

% ============================================================
\section*{5. כוח הכבידה}
% ============================================================

\subsection*{5.1 חוק הכבידה האוניברסלי של ניוטון}

\begin{tcolorbox}[defbox]
בין שני גופים נקודתיים במסות \textenglish{$m_1, m_2$} ובמרחק \textenglish{$r$} פועל כוח משיכה הדדי:
\[
  F = G\frac{m_1 m_2}{r^2},
  \qquad
  G = 6.674 \times 10^{-11}\;\text{N}\cdot\text{kg}^{-2}\cdot\text{m}^2
\]
הכוח לאורך הקו המחבר את שני הגופים, ומקיים את החוק השלישי.
\end{tcolorbox}

\subsection*{5.2 בקירוב קרוב לפני כדור הארץ}

עבור גוף במסה \textenglish{$m$} קרוב לפני כדור הארץ:
\[
  F = G\frac{m M_\oplus}{R_\oplus^2}
\]
כאשר \textenglish{$M_\oplus, R_\oplus$} הם מסת ורדיוס כדור הארץ. נסמן את הקבוע:
\[
  g \equiv \frac{GM_\oplus}{R_\oplus^2} \approx 9.8\;\text{m/s}^2
\]

\begin{tcolorbox}[resultbox]
\[
  \hvec{F}_{\text{grav}} = m\hvec{g}
  \qquad (|\hvec{g}| \approx 9.8\;\text{m/s}^2,\;\text{כלפי מטה})
\]
\end{tcolorbox}

\begin{tcolorbox}[notebox]
\textenglish{$g$} מתקבל מקבועים מסוימים של כדור הארץ. על כוכבי לכת אחרים יתקבל ערך שונה. הוא תאוצה (\textenglish{m/s$^2$}), לא כוח.
\end{tcolorbox}

% ============================================================
\section*{6. מתח בחוט}
% ============================================================

חוט יכול להפעיל כוח משיכה על גוף הקשור לקצהו. כוח זה נקרא \textbf{מתח} (\textenglish{tension}) ומסומן ב-\textenglish{$T$}.

\begin{center}
\begin{tikzpicture}[scale=1.2, >={Stealth[length=2.5mm]}]
  % Ceiling
  \draw[thick] (-1.5, 2.5) -- (1.5, 2.5);
  \fill[pattern=north east lines] (-1.5, 2.5) rectangle (1.5, 2.9);
  % Rope
  \draw[thick] (0, 2.5) -- (0, 0.8);
  % Mass
  \draw[thick, fill=blue!15] (-0.6, 0) rectangle (0.6, 0.8);
  \node at (0, 0.4) {\textenglish{$m$}};
  % Tension arrows
  \draw[->, very thick, blue!70!black] (0, 0.8) -- (0, 1.8);
  \node[blue!70!black, right] at (0.05, 1.3) {\textenglish{$\hvec{T}$}};
  % Gravity
  \draw[->, very thick, red] (0, 0.4) -- (0, -0.8);
  \node[red, right] at (0.05, -0.2) {\textenglish{$m\hvec{g}$}};
\end{tikzpicture}
\end{center}

\begin{tcolorbox}[defbox, title={חוט אידיאלי}]
חוט \textbf{חסר מסה} ובלתי-מתיח. במצב זה, המתח \textbf{אחיד לכל אורך החוט}.
\end{tcolorbox}

\textbf{נימוק:} בחוט חסר מסה, חוק שני מספר על כל קטע אינפיניטסימלי: \textenglish{$T_1 - T_2 = 0$}, ולכן \textenglish{$T_1 = T_2$} בכל נקודה.

\textbf{תוצאה מעשית:} בגלגלת חסרת חיכוך, המתח משני צדי החוט שווה.

% ============================================================
\section*{7. כוח הנורמל}
% ============================================================

\begin{tcolorbox}[defbox, title={כוח הנורמל}]
כוח שמשטח מפעיל על גוף הנוגע בו, \textbf{תמיד ניצב למשטח המגע}.
\end{tcolorbox}

\textbf{מאפיין מרכזי:} הנורמל אינו ערך קבוע מראש --- הוא \textbf{משתנה לפי הצורך} כדי למנוע חדירה לתוך המשטח. זהו \textbf{כוח אילוץ} (\textenglish{constraint force}).

\textbf{דוגמה:} קופסה על שולחן אופקי --- הנורמל יתאזן בדיוק עם מרכיב הכבידה הניצב למשטח. על שולחן בשיפוע --- הנורמל יקטן בהתאם.

% ============================================================
\section*{8. כוח החיכוך}
% ============================================================

חיכוך הוא כוח הנוצר בין שני משטחים במגע. \textbf{כיוונו תמיד מנוגד לתנועה היחסית} (או למגמת התנועה היחסית) של גוף אחד ביחס לשני.

\subsection*{8.1 חיכוך סטטי}

כאשר הגופים \textbf{אינם נעים} זה ביחס לזה, החיכוך הוא סטטי. הוא \textbf{מתאזן} עם הכוח החיצוני בדיוק כדי למנוע תנועה --- עד לגבול עליון:

\begin{tcolorbox}[resultbox]
\[
  f_s \le \mu_s N
\]
\end{tcolorbox}

\textenglish{$\mu_s$} = \textbf{מקדם החיכוך הסטטי}, תלוי בשני המשטחים. \textenglish{$N$} = כוח הנורמל ביניהם.

\subsection*{8.2 חיכוך קינטי}

כאשר החיכוך הסטטי המקסימלי מופר, הגוף מתחיל לנוע, ואז פועל \textbf{חיכוך קינטי}:

\begin{tcolorbox}[resultbox]
\[
  f_k = \mu_k N
\]
\end{tcolorbox}

\textbf{קבוע} (לא תלוי במהירות), בכיוון מנוגד למהירות היחסית.

\subsection*{8.3 גרף החיכוך כפונקציה של הכוח המופעל}

\begin{center}
\begin{tikzpicture}[scale=1.1, >={Stealth[length=2.5mm]}]
  % Axes
  \draw[->, thick] (0, 0) -- (7, 0) node[right] {\textenglish{$F_{\text{applied}}$}};
  \draw[->, thick] (0, 0) -- (0, 4.2) node[above] {\textenglish{$f$ (חיכוך)}};
  % Static (diagonal up to peak)
  \draw[very thick, blue!70!black] (0, 0) -- (3.5, 3.5);
  % Drop at threshold
  \draw[very thick, blue!70!black, dashed] (3.5, 3.5) -- (3.5, 2.5);
  % Kinetic (constant)
  \draw[very thick, blue!70!black] (3.5, 2.5) -- (6.5, 2.5);
  % Labels
  \draw[dashed, gray] (3.5, 0) -- (3.5, 3.5);
  \draw[dashed, gray] (0, 3.5) -- (3.5, 3.5);
  \draw[dashed, gray] (0, 2.5) -- (3.5, 2.5);
  \node[left] at (0, 3.5) {\textenglish{$\mu_s N$}};
  \node[left] at (0, 2.5) {\textenglish{$\mu_k N$}};
  \node[below] at (3.5, -0.05) {\small סף תזוזה};
  % Region labels
  \node[blue!70!black] at (1.5, 1.0) {\small\textit{סטטי}};
  \node[blue!70!black] at (5.0, 2.8) {\small\textit{קינטי}};
\end{tikzpicture}
\end{center}

\begin{tcolorbox}[notebox, title={פירוש הגרף}]
\begin{itemize}
  \item כל עוד \textenglish{$F_{\text{applied}} < \mu_s N$}: החיכוך הסטטי מתאים את עצמו לכוח החיצוני, הגוף נשאר במנוחה.
  \item ברגע ש-\textenglish{$F_{\text{applied}}$} עובר את \textenglish{$\mu_s N$}: הגוף מתחיל לנוע, החיכוך "קופץ למטה" לערך \textenglish{$\mu_k N$} (בדרך כלל \textenglish{$\mu_k < \mu_s$}).
  \item מכאן והלאה: החיכוך הקינטי קבוע, ללא תלות בכוח המופעל או במהירות.
\end{itemize}
\end{tcolorbox}

% ============================================================
\section*{סיכום: מבט-על}
% ============================================================

\begin{tcolorbox}[warnbox, title={נקודות מפתח לזכור}]
\begin{enumerate}
  \item \textbf{תנע} \textenglish{$\hvec{p} = m\hvec{v}$} --- הגודל ש"נשמר" במערכת סגורה.
  \item \textbf{מתקף = שינוי בתנע:} \textenglish{$\hvec{I} = \Delta\hvec{p}$}.
  \item \textbf{ניסוח כללי של חוק שני:} \textenglish{$\hvec{F} = d\hvec{p}/dt$}.
  \item \textbf{כבידה אוניברסלית:} \textenglish{$F = Gm_1 m_2/r^2$}; קרוב לפני כדה"א: \textenglish{$F = mg$}.
  \item \textbf{מתח בחוט אידיאלי} אחיד לכל אורכו.
  \item \textbf{נורמל} ניצב למשטח, מתאים את עצמו (כוח אילוץ).
  \item \textbf{חיכוך סטטי} \textenglish{$f_s \le \mu_s N$} (מתאים את עצמו עד לסף); \textbf{חיכוך קינטי} \textenglish{$f_k = \mu_k N$} (קבוע).
\end{enumerate}
\end{tcolorbox}

\end{document}